\chapter{Introduction}
\label{ch:introduction}

\section{Background and Motivation}
\label{sec:background}

Local governments in Nepal publish laws, bylaws, financial regulations, and administrative procedures as PDF documents on municipal websites. Kathmandu Metropolitan City maintains a dedicated law portal~\cite{kathmandu_law_portal}; Lalitpur and Bhaktapur publish comparable collections~\cite{lalitpur_portal,bhaktapur_portal}. These documents are written primarily in Devanagari Nepali, frequently code-switching into English for technical terms such as \emph{business registration}, \emph{property tax}, and \emph{renewal deadline}.

Citizens seeking answers to procedural questions---for example, which documents are required for a business registration in Lalitpur, or what property tax rate applies to a given valuation slab in Kathmandu---must manually search lengthy PDFs across multiple portals. Semantic retrieval systems, particularly retrieval-augmented generation (RAG) pipelines, offer a scalable alternative: documents are chunked, embedded, indexed, and queried in natural language.

However, Nepali legal PDFs pose domain-specific engineering challenges:
\begin{itemize}[leftmargin=*]
  \item \textbf{Format heterogeneity}: Some PDFs contain extractable text layers; others are scanned images requiring OCR.
  \item \textbf{Structural complexity}: Legal documents use numbered sections, annexes, schedules, and tables that affect chunk boundaries.
  \item \textbf{Linguistic code-switching}: Queries may mix Nepali and English, stressing monolingual embedding models.
  \item \textbf{Evaluation fragility}: Retrieval metrics require exact chunk-ID alignment between annotations and the indexed corpus.
\end{itemize}

This thesis addresses these challenges through \texttt{scrapktm}, a reproducible research pipeline implemented in Python.

\section{Problem Statement}
\label{sec:problem}

Given a curated set of Nepali municipal legal PDFs and a bank of realistic user queries, how can we:
\begin{enumerate}[leftmargin=*]
  \item reliably extract machine-readable text from both text-based and scanned PDFs;
  \item segment documents into retrieval-oriented chunks preserving legal structure where possible;
  \item benchmark dense semantic retrieval with reproducible metrics; and
  \item diagnose failure modes specific to Nepali legal text?
\end{enumerate}

\section{Research Objectives}
\label{sec:objectives}

The primary objectives of this work are:
\begin{itemize}[leftmargin=*]
  \item \textbf{O1}: Build a gold-standard curated corpus of 30--50 municipal and federal legal PDFs with rich metadata.
  \item \textbf{O2}: Design a hybrid extraction router that selects PyMuPDF or Surya OCR based on document format type.
  \item \textbf{O3}: Implement semantic chunking with heading detection and fixed-size fallback (512 tokens, 128 overlap).
  \item \textbf{O4}: Construct an 11-query evaluation bank targeting realistic failure modes (list extraction, table lookup, cross-references).
  \item \textbf{O5}: Compare \texttt{all-MiniLM-L6-v2} and \texttt{BAAI/bge-m3} using Recall@5 and MRR on a local FAISS index.
\end{itemize}

\section{Contributions}
\label{sec:contributions}

This thesis makes the following contributions:
\begin{itemize}[leftmargin=*]
  \item A curated 50-document Nepali legal corpus spanning four sources and four thematic categories.
  \item A modular hybrid extraction pipeline with OCR accuracy tracking and structured JSONL outputs.
  \item A code-switched research query bank with annotated ground-truth locations.
  \item A ground-truth mapping utility resolving chunk-ID schema mismatches between annotation and corpus.
  \item An open, offline-capable embedding benchmark using FAISS and sentence-transformers.
\end{itemize}

\section{Thesis Organization}
\label{sec:organization}

\Cref{ch:related} reviews related work on embeddings, OCR, and legal IR.
\Cref{ch:methodology} describes the end-to-end pipeline architecture.
\Cref{ch:corpus} details corpus curation and statistics.
\Cref{ch:evaluation} presents the query bank, metrics, and benchmark protocol.
\Cref{ch:results} reports experimental results and failure analysis.
\Cref{ch:discussion} discusses implications and limitations.
\Cref{ch:conclusion} concludes and outlines future work.
